\chapter{The dashboard, view by view}\label{ch:dashboard}

The engine serves its own web page at \code{/dashboard/}, and the standalone proxy serves the same page at \code{:8080/dashboard/}. It answers what a log cannot answer at a glance: what the engine is doing right now, for whom, and how fast. It is written in TypeScript with Lit components, its charts are hand-drawn SVG, it loads nothing from a CDN, and the server serves it from \code{dashboard/dist/} with no Node on the machine that serves it. The parts every view loads are about 45~KB of gzip. Figure~\ref{fig:dashboard-map} shows the seven views and what each reads.

\begin{figure}[htbp]
\centering
% The dashboard's views and the server routes each reads.
\begin{tikzpicture}[
    tab/.style={draw=cobalt, fill=cobaltlight, rounded corners=2pt, font=\sffamily\small,
      minimum width=24mm, minimum height=7mm},
    src/.style={font=\ttfamily\scriptsize, text=slate, anchor=west, align=left},
    node distance=2.4mm]
  \node[tab] (usage) {Usage};
  \node[tab, below=of usage] (perf) {Performance};
  \node[tab, below=of perf] (dev) {Dev};
  \node[tab, below=of dev] (play) {Playground};
  \node[tab, below=of play, fill=signallight, draw=signal] (cond) {Conductor};
  \node[tab, below=of cond] (back) {Backends};
  \node[tab, below=of back] (sys) {System};
  \node[src] at ([xshift=4mm]usage.east) {/v1/dashboard/summary, /usage \quad a year of tokens, cost by backend and user};
  \node[src] at ([xshift=4mm]perf.east) {/v1/dashboard/live (SSE), /requests, /metrics, /parallel};
  \node[src] at ([xshift=4mm]dev.east) {/v1/dashboard/logs \quad the live log, a request's detail};
  \node[src] at ([xshift=4mm]play.east) {/v1/chat/completions \quad the OpenAI API itself};
  \node[src] at ([xshift=4mm]cond.east) {/v1/conductor/* \quad decisions, documents, sessions, memory};
  \node[src] at ([xshift=4mm]back.east) {/v1/config/* \quad backends, models, prices, rules, users};
  \node[src] at ([xshift=4mm]sys.east) {/v1/dashboard/system \quad settings, memory plan, hosts};
\end{tikzpicture}

\caption[The dashboard's views]{The dashboard's seven views, in the order of its navigation rail, and the server routes each reads. Every view only reads, except where an edit makes a new version of a document or a configuration row.}\label{fig:dashboard-map}
\end{figure}

\section{Performance: the live panel}\index{Live panel}

The Performance view has two pages. Engine shows the engine's own requests; Parallel shows every request on every backend (Section~\ref{sec:parallel}).

\screenshot{dashboard-performance}{The Performance view, Engine page: the Live panel with a request writing, the decode and prefill rates, and the last requests. Demo data from the dashboard's mock server (\texttt{npm run preview:mock}), not a measurement.}{fig:shot-performance}

The Live panel on top streams \code{GET /v1/dashboard/live} over server-sent events: four events a second while a request runs, one otherwise. A line at the top says what the engine is doing now, in the server's own words:

\begin{itemize}
\item \texttt{Prefilling 36,864 of 48,210 (76\,\%)}, with the prefill rate and the time left;
\item \enquote{Thinking} or \enquote{Writing}, with the decode rate and the tokens so far;
\item \enquote{Calling tool write\_file}, with the size of the arguments written so far;
\item \enquote{Waiting for client: running tool bash}, between two turns of an agent;
\item \enquote{Idle} or \enquote{Draining}.
\end{itemize}

\begin{figure}[htbp]
\centering
% The Performance view's Live panel as an instrument readout. The numbered marks are
% explained in the caption. The figures shown are examples, not measurements.
\begin{tikzpicture}[x=1mm, y=1mm, font=\sffamily\scriptsize,
    seg/.style={minimum height=2.6mm, inner sep=0pt, outer sep=0pt},
    callout/.style={circle, draw=cobalt, fill=white, text=cobalt, font=\sffamily\bfseries\scriptsize,
      inner sep=1pt, minimum size=4.2mm}]
  % the panel
  \node[panel, anchor=north west, minimum width=136mm, minimum height=76mm] (p) at (0,0) {};
  % the state line
  \node[anchor=west, font=\sffamily\small, text=ink] at (4,-5) {\textbf{Thinking} \quad 49.1 tok/s \quad 4,812 tokens so far};
  \node[anchor=east, text=slate] at (126,-5) {streaming $\cdot$ 4/s};
  \node[callout] at (132,-5) {1};
  \draw[rule] (4,-8.5) -- (132,-8.5);
  % the two big figures with sparklines
  \node[anchor=north west, readout] at (4,-10) {49.1};
  \node[anchor=north west, unit] at (4,-17) {decode tok/s};
  \draw[cobalt, thick] (26,-17) -- (30,-14) -- (34,-15) -- (38,-12) -- (42,-13.5) -- (46,-11) -- (50,-12) -- (54,-11.5) -- (58,-13) -- (62,-11.5);
  \node[anchor=north west, unit] at (26,-17.5) {last 5 minutes};
  \node[anchor=north west, readout] at (70,-10) {742};
  \node[anchor=north west, unit] at (70,-17) {prefill tok/s};
  \draw[cobalt, thick] (90,-16) -- (94,-12) -- (98,-12.5) -- (102,-11) -- (106,-16) -- (110,-16) -- (114,-11.5) -- (118,-12) -- (122,-13) -- (126,-12.5);
  \node[anchor=north west, unit] at (90,-17.5) {last 5 minutes};
  \node[callout] at (132,-13) {2};
  % tokens per round and ms per round
  \node[anchor=west, text=ink] at (4,-23.5) {4.6 tokens a round \quad 93 ms a round \quad {\color{slate} over the last second}};
  \node[callout] at (132,-23.5) {3};
  \draw[rule] (4,-26.5) -- (132,-26.5);
  % one request in flight: the lane
  \node[anchor=west, text=ink] at (4,-30.5) {\textbf{chatcmpl-9c2e} \quad opencode \quad coding \quad thinking 3,900 + content 912};
  \node[anchor=east, text=slate] at (126,-30.5) {TTFT 1.8 s};
  \node[callout] at (132,-32.5) {4};
  % the timeline strip, one segment per state
  \node[seg, fill=slate!35, minimum width=14mm, anchor=west] at (4,-35) {};
  \node[seg, fill=cobalt!45, minimum width=54mm, anchor=west] at (18,-35) {};
  \node[seg, fill=cobalt, minimum width=26mm, anchor=west] at (72,-35) {};
  \node[seg, fill=signal!70, minimum width=10mm, anchor=west] at (98,-35) {};
  \node[unit, anchor=north] at (11,-36.5) {prefill};
  \node[unit, anchor=north] at (45,-36.5) {thinking};
  \node[unit, anchor=north] at (85,-36.5) {writing};
  \node[unit, anchor=north] at (103,-36.5) {tool call};
  \node[anchor=west, text=slate] at (4,-42) {› writing render.js, 211 lines so far \qquad continues chatcmpl-7f3a after 8.4 s, client ran bash};
  \node[callout] at (132,-42) {5};
  \draw[rule] (4,-45.5) -- (132,-45.5);
  % the last requests
  \node[anchor=west, unit] at (4,-48.5) {last 20 requests};
  \node[callout] at (132,-56) {6};
  \node[anchor=west, text=ink] at (4,-52.5) {\textcolor{moss}{\rule{1.6mm}{1.6mm}}\; stop after 1,204 tokens, 48.7 tok/s, cache \texttt{end}};
  \node[anchor=west, text=ink] at (4,-56.5) {\textcolor{signal}{\rule{1.6mm}{1.6mm}}\; abandoned by the client after 73 s of silent prefill, 0 tokens sent};
  \node[anchor=west, text=ink] at (4,-60.5) {\textcolor{moss}{\rule{1.6mm}{1.6mm}}\; tool\_calls after 388 tokens, 51.2 tok/s, side lane};
  \node[anchor=west, text=ink] at (4,-64.5) {\textcolor{signal!60!yellow}{\rule{1.6mm}{1.6mm}}\; refused: queue full (503, Retry-After 5)};
  \node[anchor=west, unit] at (4,-70.5) {red: abandoned, failed, timed out, cancelled \quad amber: refused, rejected};
\end{tikzpicture}

\caption[The Live panel as an instrument readout]{The Live panel, read as an instrument. (1) The state line, in the server's own words, and the panel's report on its own stream. (2) The two big figures with 5-minute sparklines; each leaves the request's warm-up out and reads a dash when nothing runs. (3) Tokens a round and milliseconds a round, the two factors of the engine's speed. (4) One lane per request in flight: state, client, profile, tokens split into thinking and content, time to first token, and a strip with one segment per state. (5) The newest status line, and what the agent did between two turns. (6) How each of the last 20 requests ended, in one sentence; a row opens its detail in the Dev view. The figures shown are examples, not measurements.}\label{fig:live-panel}
\end{figure}

When the client's socket closes while the engine still works for it, the line says so in red: \enquote{client disconnected 12 s ago}. Under the two big figures (decode and prefill tok/s, each with a 5-minute sparkline) the panel shows tokens per round and milliseconds per round over the last second. Each request in flight gets a row: its state, its client, tokens split into thinking and content, time to first token, the rates, the newest status line, and a timeline strip with one segment per state. When an agent sends its next turn, the row says which request it continues and what the client did in between (\enquote{continues chatcmpl-7f3a after 8.4 s, client ran bash}).

The live rates leave each request's warm-up out (the first 50 decode tokens and the first prefill chunk) and show a dash when nothing runs, so a figure is always a rate of work being done now. A list of the last 20 requests says how each one ended in one sentence (\enquote{abandoned by the client after 73 s of silent prefill, 0 tokens sent}); abandoned, failed, timed-out and cancelled requests get a red mark, refused and rejected ones an amber one. The panel also reports on its own stream (\enquote{streaming · 4/s}, \enquote{no update for 3.0 s}, \enquote{reconnecting since 4 s}), so a silent page is never mistaken for an idle engine.

Every state and every sentence comes from the server (\code{server/activity.py}); the page only formats numbers. The engine's loop keeps the state it needs anyway, and a sampler reads it once per tick, so no statement runs per token for the dashboard's sake. A tick costs 410 to 470~\textmu s of CPU at 4~Hz (\measured), and with the activity on and a stream open, the release gate found the text byte-identical and the round time unchanged within noise. \code{--live-activity off} turns the activity fields off.

\section{Performance: the history}

Below the Live panel: the time to first token, tokens per round and acceptance over the last 5 minutes, then speed history from the usage ledger. The Recent requests chart plots decode tok/s against prompt length, one dot per request. A replay (a response-cache hit, a summary written ahead) decoded nothing and is never a dot there, nor in any average, percentile or live rate; a switch lists the replays with their time to first token. Side-lane requests are not plotted either, because a short plain decode beside a turn says nothing about the engine's decode rate. The y axis ends at twice the 95th percentile, and a point above it is a triangle at the top edge with its value in the tooltip. Each request's cache source (\code{end}, \code{resident}, \code{disk}, \code{spill}, \code{continuation} and the others) is shown with it, so a slow first token can be read at once.

\section{Usage}\index{Usage view}

A year of tokens as a heatmap, with totals for today, 7 days, 30 days and the year, split four ways: input, cached input, output, reasoning. Filters by model, client, backend (\code{engine} is the engine's own rows) and user, and the cost of priced models. The view reads the usage ledger, which keeps at least 400 days and never stores prompt or answer text (Section~\ref{sec:ledger}). TandemLLM's own calls to itself (the router's question to the engine, client kind \code{internal}) are left out of the totals, the filters and Dev's recent requests unless the \enquote{show internal} switch is on; each names the request it served. In the Live panel the cached tokens never read above the client's prompt, and TandemLLM's notes show as their own figure, \enquote{+5,235 hidden}.

\screenshot{dashboard-usage}{The Usage view: a year of tokens, requests and cost, filtered by model, client, backend and user. Demo data from the dashboard's mock server (\texttt{npm run preview:mock}), not a measurement.}{fig:shot-usage}

\section{Dev}

The live server log, filterable, with a detail view for each request: its timings, its cache source, its tokens, the conductor's decisions for it and the reason it ended. A memory-gate refusal is an error line here with the allocated and reserved figures.

\screenshot{dashboard-dev}{The Dev view: the live log, which never holds prompt or answer text, and the recent requests. Demo data from the dashboard's mock server (\texttt{npm run preview:mock}), not a measurement.}{fig:shot-dev}

\section{Playground}

A streaming chat against the server's own OpenAI API, with tools and parameters, and the figures of each answer. It uses the same routes any client uses, so it works with the login on or off.

\screenshot{dashboard-playground}{The Playground before the first message: the chat on the left, the system prompt and the sampling parameters on the right, the engine's own figures for each turn below. Demo data from the dashboard's mock server (\texttt{npm run preview:mock}), not a measurement.}{fig:shot-playground}

\section{Conductor}\index{Conductor view}

The Conductor view shows how \code{TandemLLM} decides, and edits every knob of it. Every edit is checked first, shows its diff, and becomes a new version with a note; nothing asks for a password. Its screens:

\begin{table}[htbp]
\centering\small
\caption{The Conductor view's screens.}\label{tab:conductor-screens}
\begin{tabularx}{\linewidth}{@{}L{.14\linewidth}Y@{}}
\toprule
\hd{screen} & \hd{what it shows and edits} \\
\midrule
Overview & the live decision feed: each turn's class, confidence, effort, note and the monitor's events; tiles over the last 24 hours \\
Routing & the task classes and their policy (effort, effort by size, packages, rules), the classifier stages and thresholds, stickiness and the structural rules, a test box that classifies a text you type \\
Monitor & the session monitor's journal and rules: the tool checks' floors and project rules, the reasoning arms, the empty-turn nudges, the session tasks with their texts and patterns, and what they hold per session right now \\
Profiles & each profile's switches (mode, decisions, tasks and caps, thinking, what the user sees, outcomes, library, memory, shadow), the roles with their contracts, and the residency plan \\
Packages & every package, task skill, test skill, card and library skill, with tokens against limits and recent use; the editor; the Library tab; the Proposals tab \\
Deciders & the decision models: state, limits, latency, load and unload through the residency planner, the questions each answers, Laya for scripts \\
Sessions & the sessions, each a timeline of its turns and every hidden text the model saw, with an export \\
Memory & three tabs: the session transcripts kept on the machine, the memory notes (Chapter~\ref{ch:memory}), and the disk tier's state and settings \\
Compare & the shadow precision of each task and question, per backend, from the labels a person gave decisions, and the evaluation reports \\
\bottomrule
\end{tabularx}
\end{table}

\screenshot{dashboard-conductor-overview}{The Conductor view, Overview: decisions and the session monitor's acts over 24 hours, the classifier's latency, the roles online and the task classes. Demo data from the dashboard's mock server (\texttt{npm run preview:mock}), not a measurement.}{fig:shot-conductor}

\screenshot{dashboard-conductor-sessions}{The Conductor view, Sessions: each session's client, profile, current class, turns and the hidden notes it carries. Demo data from the dashboard's mock server (\texttt{npm run preview:mock}), not a measurement.}{fig:shot-sessions}

\screenshot{dashboard-conductor-memory}{The Conductor view, Memory, the Notes tab: the live notes by scope, with kind, status, writer and use. Demo data from the dashboard's mock server (\texttt{npm run preview:mock}), not a measurement.}{fig:shot-memory}

\section{Backends}\index{Backends view}

On an engine without \code{--backends} the System answer says \code{proxy.mode: off}, and the view says the proxy is off without asking the configuration routes. The dashboard has no token dialog and no sign-in screen (Section~\ref{sec:dash-login}).

\screenshot{dashboard-backends}{The Backends view with a vLLM, an SGLang and a DeepSeek backend, and the line for this engine's own settings. Demo data from the dashboard's mock server (\texttt{npm run preview:mock}), not a measurement.}{fig:shot-backends}

The proxy's registry: each backend's health, models, requests in flight and cost today. Add one after a test connection, edit, delete. The key field is write-only and shows the key's last 4 characters; a field a \code{TANDEM_*} variable sets is locked and names it. The models have public ids, prices (USD per million tokens, the release's list prices one click away) and the conductor's feature switches, each with what it resolves to and what it costs on that backend. The background memory writer's switches are on a backend's form. The ordered routing rules and the users with their keys, each key shown once, complete it. With no backend the page opens on Setup; on an engine without \code{--backends}\index{backends@\texttt{-\kern0pt-backends}} it says how to turn the proxy on. A vLLM or SGLang backend's card shows its engine's own figures and whether its plugin reports (\enquote{L2 plugin} or \enquote{L1 only}, Chapter~\ref{ch:hosts}).

\section{System}\index{System view}

The engine's settings (its command line and every \code{QWEN38_*} and \code{QSE_*} variable, the home directory written as \code{~}, every token-like value redacted) and the machine's state: memory, GPU, caches, with 30 minutes of history kept in the browser. The memory plan shows the engine's own split (weights, KV pool, drafter caches, graphs, each read off the object that holds it), \code{MemAvailable}, every role with its claim and whether it fits now, the limits with what each holds, and the last refusal. Below them are the hosts: this machine with the engine's process and every backend on it, and each remote backend's machine, down to the GPU memory of its processes (Chapter~\ref{ch:hosts}).

The engine card carries one neutral line, the speed check: \enquote{passed for 0.4.0 · <date>} when the build has a speed guard report (the release ships its own, \code{release/speed-guard.json}), \enquote{not run for this build} otherwise; the numbers are in its tooltip and on the Conductor's Compare screen. It is no warning: the speed guard is the release process's check, not something to act on in the dashboard.

\screenshot{dashboard-system}{The System view: the engine, the memory, the GPU, the queue, the caches and the machine. Demo data from the dashboard's mock server (\texttt{npm run preview:mock}), not a measurement.}{fig:shot-system}

\section{The contract}\index{dashboard contract}

\code{docs/contract/dashboard-v1/} holds a JSON Schema for every JSON response of \code{/v1/dashboard/*} and a valid example of each; \code{docs/contract/conductor-v1/} does the same for the conductor's routes. Both sides test against them: the server's tests validate real responses, and the dashboard's tests validate the examples and every response of its mock server. A change to a field bumps \code{contract_version}, and the contract's README lists every change.

\section{Building and testing the page}

\begin{shell}
cd dashboard
npm ci
npm run dev:mock        # the app against a mock engine
npm test                # unit, contract and build tests
npm run build           # type-check and build dist/
npx playwright test     # end-to-end, Chromium and WebKit, desktop and phone width
\end{shell}

The mock engine plays scenarios for the Live panel (an agent turn with a chunked prefill and a tool call, a client that leaves during a prefill, one request per stop reason) and failure modes (an expired session, server errors, dropped streams). The end-to-end suites also run against \code{server/app.py --fake-engine}, which serves the real API with a scripted engine and no GPU. Every view works at a phone's width of 390 pixels: the navigation becomes a bottom bar with short labels, and lanes and tables become cards.
